\begin{ExerciseBlock}
\ExerciseHeading{02}{Fermatsches Prinzip: Reflexion und Brechung}{geometrical_and_wave_optics_fermat_principle}
Nach dem Fermatschen Prinzip ist die Laufzeit des Lichts entlang des tatsächlichen Strahlwegs stationär. In den beiden unten betrachteten Anordnungen wird die Laufzeit minimiert.
\begin{enumerate}
\item Bei der Reflexion läuft der Strahl von $A$ über den Reflexionspunkt $P$ nach $B$. Berechnen Sie die Laufzeit $\tau(x_P)$. Zeigen Sie, dass sie genau dann minimal ist, wenn Einfalls- und Ausfallswinkel gleich sind.
\begin{center}
\CachedTikzFigure{content/resources/chapters/geometrical_and_wave_optics/tikz/fermat_reflection.tex}
\end{center}
\item Bei der Brechung läuft der Strahl von $A$ über den Punkt $P$ an der Grenzfläche nach $B$. Berechnen Sie erneut $\tau(x_P)$ und zeigen Sie, dass sie minimal ist, wenn $n_1\sin\vartheta_1=n_2\sin\vartheta_2$ gilt.
\begin{center}
\CachedTikzFigure{content/resources/chapters/geometrical_and_wave_optics/tikz/fermat_refraction.tex}
\end{center}
\end{enumerate}
\end{ExerciseBlock}
\UseExerciseSolution{geometrical_and_wave_optics_fermat_principle}
